% ==============================================================================
% PHỤ LỤC A: MÃ NGUỒN CHƯƠNG TRÌNH
% ==============================================================================

\chapter{MÃ NGUỒN CHƯƠNG TRÌNH}
\label{app:sourcecode}

\section{Liên kết kho mã nguồn (Repository)}

Toàn bộ mã nguồn, dữ liệu mẫu và tài liệu hướng dẫn chi tiết của dự án được lưu trữ công khai tại kho chứa mã nguồn GitHub:
\begin{itemize}[leftmargin=1.5cm]
    \item \textbf{Địa chỉ repository:} \url{https://github.com/ictu-fit/ai-app-template}
    \item \textbf{Nhánh chính (Main branch):} \texttt{main}
    \item \textbf{Giấy phép mã nguồn:} MIT License
\end{itemize}

\section{Cấu trúc thư mục dự án phần mềm}
 
Dưới đây là sơ đồ tổ chức thư mục của dự án phần mềm ứng dụng AI:
\begin{lstlisting}[language={}, caption={Cấu trúc cây thư mục mã nguồn phần mềm ứng dụng AI}, label={lst:tree}]
AI_Application/
|-- data/                 # Thư mục lưu trữ tập dữ liệu (train, val, test)
|-- models/               # Định nghĩa các kiến trúc mạng nơ-ron
|   |-- __init__.py
|   |-- backbone.py       # Module trích xuất đặc trưng
|   `-- classifier.py     # Module phân loại
|-- weights/              # Tệp trọng số mô hình đã huấn luyện (.pt, .onnx)
|-- utils/                # Các hàm tiện ích tiền xử lý dữ liệu và vẽ đồ thị
|   |-- dataset.py
|   `-- metrics.py
|-- ui/                   # Giao diện người dùng (Streamlit / PyQt)
|   `-- app.py
|-- train.py              # Script chạy quá trình huấn luyện
|-- evaluate.py           # Script đánh giá mô hình trên tập test
|-- requirements.txt      # Danh sách các thư viện phụ thuộc
`-- README.md             # Hướng dẫn tổng quan về dự án
\end{lstlisting}
